\section{Objectifs}
\begin{itemize}
\item Calculer des aires par intégration.
\item Relier intégrale et primitive.
\item Utiliser des approximations par rectangles.
\end{itemize}

\section{1. Primitive}
$F$ est une primitive de $f$ si
\[
F'=f.
\]

\section{2. Intégrale}
Pour $f$ continue :
\[
\int_a^bf(x)\,dx=F(b)-F(a).
\]

\section{3. Aire positive}
Si $f\ge0$ sur $[a,b]$, l'intégrale représente l'aire sous la courbe.

\section{4. Aire entre deux courbes}
Si $f\ge g$ :
\[
\mathcal A=\int_a^b(f-g).
\]

\section{5. Méthode des rectangles}
On découpe l'intervalle en sous-intervalles et on additionne les aires de rectangles.

Cette méthode donne une approximation qui peut être améliorée en raffinant le découpage.

\section{6. Valeur moyenne}
\[
\mu=\frac1{b-a}\int_a^bf(x)\,dx.
\]

\section{Erreurs fréquentes}
\begin{itemize}
\item Confondre aire géométrique et intégrale signée.
\item Oublier les bornes.
\item Utiliser une primitive incorrecte.
\end{itemize}

\section{À retenir}
L'intégrale donne une mesure exacte ou approchée d'aire et relie géométrie et analyse.
